%%% ==========================================================================
%%% CNN-Based Semiconductor Wafer Defect Detection
%%% Historical report source retained for reference
%%% ==========================================================================
\documentclass[conference]{IEEEtran}

% ---------- packages --------------------------------------------------------
\usepackage{amsmath,amssymb}
\usepackage{graphicx}
\usepackage{booktabs}
\usepackage[hidelinks]{hyperref}
\usepackage{float}
\usepackage{cite}
\usepackage{array}
\usepackage{multirow}
\usepackage{url}

% ============================================================================
\begin{document}

% ---------- title -----------------------------------------------------------
\title{CNN-Based Semiconductor Wafer Defect Detection\\Using the WM-811K Dataset}

\author{
\IEEEauthorblockN{Anindita Paul, Brandon Pyle, Anand Rajan, Brett Rettura}
\IEEEauthorblockA{Historical report source retained for reference\\
Current verified defense narrative: \texttt{docs/DEFENSE\_PACKET.md}}
}

\maketitle

% ---------- abstract --------------------------------------------------------
\begin{abstract}
Semiconductor manufacturing demands rapid, reliable detection of wafer-level
defect patterns to maintain high yield and reduce production cost. In this
repository centers a modular software artifact for wafer-defect classification
and analysis using the WM-811K dataset. The codebase includes a custom CNN,
transfer-learning baselines, explainable inference through Grad-CAM,
uncertainty estimation, anomaly and OOD analysis, federated-learning utilities,
synthetic augmentation components, and a versioned model-registry path. This
report source is retained as a historical narrative document; the current
verified defense-facing status, validation evidence, and limitations are
summarized in \texttt{docs/DEFENSE\_PACKET.md} and
\texttt{docs/FINAL\_STATUS\_REPORT.md}.
\end{abstract}

\begin{quote}
\small
\textbf{Note:} This document is preserved as a historical project report source.
For committee use, treat \texttt{docs/DEFENSE\_PACKET.md} and
\texttt{docs/FINAL\_STATUS\_REPORT.md} as the authoritative summary of the
current verified repository state.
\end{quote}

\begin{IEEEkeywords}
wafer defect detection, convolutional neural network, WM-811K, transfer
learning, class imbalance, Grad-CAM
\end{IEEEkeywords}

% ============================================================================
\section{Introduction}
\label{sec:introduction}
% ============================================================================

The cornerstone of the current economy is the semiconductor industry. The
projected 2026 annual sales value of this industry is US\$975
billion~\cite{sia2024}, driven by intense demand for AI infrastructure. As
device size has continued to shrink in accordance with Moore's Law, with
current nodes at 3\,nm and below, the complexity of the manufacturing
processes has increased exponentially. A single semiconductor wafer can contain
thousands of individual chips, each worth hundreds or thousands of dollars.

Semiconductor manufacturing combines repetitive, complex, additive, and
subtractive processes to assemble integrated circuit (IC) chips on a silicon
wafer, making it among the most complex manufacturing
processes~\cite{kahng2021}. It can be classified as (a)~the \emph{front-end
process} and (b)~the \emph{back-end process}.

In the front-end process, multiple circuit layers are deposited via the
following key steps. First, \textbf{oxidation and deposition}: layers of
insulating material such as SiO$_2$ are deposited via oxidation, and using
chemical vapor deposition (CVD) or physical vapor deposition (PVD), metals are
deposited on the wafer surface to form conductive layers. Second, precise
patterns are defined by means of \textbf{photolithography and etching}.
Third, \textbf{ion implantation} or diffusion introduces impurities to alter
the electrical properties, such as conductivity, to form transistors or other
devices.

The effectiveness of the front-end process is assessed through wafer testing,
a back-end process. In this wafer-test step the number of functional chips is
counted and the ratio of functional chips to total chips is defined as the
\emph{yield}---the key indicator of semiconductor manufacturing process
quality~\cite{weng2010}. Due to technological scaling, the minimum feature
size is now below 3\,nm. As a result, a single microscopic defect that is
invisible to the naked eye can cause an entire die to fail. These defects
manifest as distinct spatial patterns on the wafer and can arise from hundreds
of complex process steps such as lithography, etching, deposition, and
polishing. Identification and classification of defect patterns in wafer maps
can provide engineers with useful information to detect the root causes of
failures in wafer manufacturing, thereby improving wafer production
yield~\cite{chen2023}.

% ============================================================================
\section{Problem Statement}
\label{sec:problem}
% ============================================================================

The semiconductor manufacturing industry faces significant challenges in
accurately and efficiently detecting and classifying wafer defects during
production. Current manual inspection methods are time-consuming, subjective,
and error-prone, leading to inconsistent defect classification, delayed
production cycles, and substantial yield losses.

This project develops an automated detection and classification system using
three CNN architectures---a custom CNN, ResNet-18~\cite{he2016}, and
EfficientNet-B0~\cite{tan2019}---on the WM-811K dataset~\cite{wu2015}, the
largest publicly available dataset with 811K wafer maps across nine classes.
Our deep-learning models mitigate class imbalance and noise to achieve
accurate pattern recognition, treating wafer maps as images. The successful
implementation of this automated system aims to enhance manufacturing yield by
eliminating human variability, providing uniform defect classification 24/7,
and enabling early fault detection to prevent defect propagation.

% ============================================================================
\section{Challenges}
\label{sec:challenges}
% ============================================================================

Several challenges arise in applying deep learning to wafer-map defect
detection:

\begin{enumerate}
    \item \textbf{Severe class imbalance.} The WM-811K dataset is dominated by
          the ``none'' class ($\sim$85\% of labeled samples), while the rarest
          defect class (Near-full) contains only 149 examples. Standard
          classifiers trained without mitigation strategies will overwhelmingly
          predict the majority class~\cite{he2009}.

    \item \textbf{Computational constraints.} Training deep CNNs on a dataset
          of 172{,}950 wafer maps is computationally expensive. In this study
          we are limited to CPU-only hardware, necessitating architectural
          compromises such as pre-resizing to $96 \times 96$ and capping the
          number of training samples per epoch.

    \item \textbf{Domain gap for transfer learning.} Pretrained ImageNet
          models have learned features from natural photographs, whereas wafer
          maps are synthetic, low-resolution binary-like images. The
          transferability of such features to the wafer-map domain is not
          guaranteed~\cite{yosinski2014}.

    \item \textbf{Defect pattern diversity.} The nine defect classes exhibit
          highly variable spatial morphologies. Some patterns (e.g., Edge-Ring)
          are visually distinctive, while others (e.g., Loc vs.\ Random) can be
          difficult to distinguish even for experienced
          engineers~\cite{chen2023}.

    \item \textbf{Legacy data format.} The WM-811K dataset is stored as a
          Python~2 pickle file, requiring encoding workarounds and module-path
          remapping for modern Python~3 environments.
\end{enumerate}

% ============================================================================
\section{Related Work}
\label{sec:related}
% ============================================================================

Wu et al.~\cite{wu2015} introduced the WM-811K dataset and proposed a
similarity-based ranking approach using density-based spatial features for
root-cause analysis. Nakazawa and Kulkarni~\cite{nakazawa2018} subsequently
applied a CNN architecture to WM-811K, coupling defect classification with
image retrieval. Yu et al.~\cite{yu2019} proposed a hybrid framework
combining a densely connected CNN with deep forest, leveraging transfer
learning for minority classes. Saqlain et al.~\cite{saqlain2020} introduced a
voting ensemble classifier that aggregated predictions from multiple base
learners. Kahng and Kim~\cite{kahng2021} explored self-supervised
representation learning for wafer bin map classification, demonstrating that
pre-text tasks tailored to the spatial structure of wafer maps can yield
informative features without labeled data. Kim et al.~\cite{kim2022} proposed
a single-shot detector for identifying mixed-type defect patterns. Chen et
al.~\cite{chen2023} developed an attention-based method for wafer-map defect
pattern detection.

Foundational deep-learning work underpinning our approach includes
ResNet~\cite{he2016} with skip connections for training very deep networks,
EfficientNet~\cite{tan2019} with compound scaling, batch
normalization~\cite{ioffe2015}, dropout~\cite{srivastava2014}, and
Grad-CAM~\cite{selvaraju2017} for visual explanations. Transfer learning from
ImageNet~\cite{pan2010,yosinski2014} has proven effective across industrial
imaging tasks~\cite{weimer2016}.

Despite substantial progress, several gaps remain: many prior studies report
only overall accuracy without per-class metrics for imbalanced
data~\cite{he2009,chawla2002}; few provide visual explanations of model
decisions; and comparisons across modern efficient architectures are
underrepresented. Our work addresses these gaps by systematically comparing
three CNN architectures under identical conditions, employing multiple
class-imbalance strategies, and providing Grad-CAM visualizations.

% ============================================================================
\section{Data Description}
\label{sec:data}
% ============================================================================

\subsection{The WM-811K Dataset}

The WM-811K dataset~\cite{wu2015} is distributed as a serialized Python
pickle file (\texttt{LSWMD.pkl}) and contains 811{,}457 wafer map records
collected from real semiconductor production lines. Each record includes a
two-dimensional array where each element takes one of three values: 0
(background), 1 (normal die), or 2 (defective die), along with metadata such
as lot name, wafer index, and failure type label.

Only the subset of wafers with an assigned failure-type label is used for
supervised classification. After filtering, approximately 172{,}950 labeled
wafer maps remain, categorized into the following nine classes of systematic
defect patterns~\cite{chen2023,weng2010}:

\begin{itemize}
    \item \textbf{Center}: concentrated at the wafer center, caused by
          failure in thin-film deposition~\cite{chen2021}.
    \item \textbf{Donut}: ring-shaped pattern from photoresist redeposition or
          chemical dispense issues.
    \item \textbf{Edge-Ring}: along the wafer edge from abnormal etching.
    \item \textbf{Loc (Local)}: small cluster from non-uniform developer
          application or pH changes.
    \item \textbf{Edge-Loc}: localized defect along the wafer edge from
          edge-sensitive processes.
    \item \textbf{Near-full}: high-density defects covering most of the wafer
          from critical process failure.
    \item \textbf{Scratch}: linear patterns from physical damage or debris.
    \item \textbf{Random}: irregularly distributed defects from cleanroom
          contaminants.
    \item \textbf{None}: no systematic spatial pattern detected.
\end{itemize}

\subsection{Class Distribution}

Table~\ref{tab:class_distribution} shows the severe class imbalance: the
``none'' class accounts for over 85\% of labeled samples, while Near-full
contains only 149 examples ($\sim$990:1 ratio).

\begin{table}[H]
\centering
\caption{Class distribution of labeled wafer maps.}
\label{tab:class_distribution}
\begin{tabular}{@{}lrr@{}}
\toprule
\textbf{Defect Class} & \textbf{Count} & \textbf{\%} \\
\midrule
None        & 147{,}431 & 85.24 \\
Edge-Ring   &   9{,}680 &  5.60 \\
Edge-Loc    &   5{,}189 &  3.00 \\
Center      &   4{,}294 &  2.48 \\
Loc         &   3{,}593 &  2.08 \\
Scratch     &   1{,}193 &  0.69 \\
Random      &     866   &  0.50 \\
Donut       &     555   &  0.32 \\
Near-full   &     149   &  0.09 \\
\midrule
\textbf{Total} & \textbf{172{,}950} & \textbf{100.00} \\
\bottomrule
\end{tabular}
\end{table}

\subsection{Preprocessing}

Each wafer map is preprocessed as follows: (1)~resized to $96 \times 96$
pixels using nearest-neighbor interpolation to preserve the discrete die
values; (2)~pixel values normalized to $[0, 1]$ by dividing by~2; and
(3)~the single-channel image is replicated across three channels to produce a
$3 \times 96 \times 96$ tensor compatible with pretrained ImageNet models.

% ============================================================================
\section{Deep Learning Methodology}
\label{sec:methodology}
% ============================================================================

This is a \textbf{supervised multi-class classification} problem: given a
wafer map image, the model predicts which of nine defect classes it belongs
to. We employ convolutional neural networks (CNNs), the standard approach for
image classification since the landmark work of Krizhevsky et
al.~\cite{krizhevsky2012}.

\subsection{Data Pipeline}

The labeled dataset is split into training (70\%), validation (15\%), and test
(15\%) subsets via stratified sampling. During training, a
\texttt{WeightedRandomSampler} assigns each sample a weight equal to the
inverse frequency of its class:
\begin{equation}
w_i^{\text{sample}} = \frac{1}{n_{c(i)}}
\label{eq:sample_weight}
\end{equation}
where $n_{c(i)}$ is the training count for class $c(i)$ of sample $i$.
Data augmentation includes random horizontal/vertical flips, rotation
($\pm 15^\circ$), and affine translation (up to 5\%).

The loss function is class-weighted cross-entropy:
\begin{equation}
\mathcal{L} = -\sum_{c=1}^{9} w_c \, y_c \, \log(\hat{p}_c)
\label{eq:weighted_ce}
\end{equation}
where $w_c$ is inversely proportional to class frequency, further penalizing
misclassification of minority classes.

\subsection{Model Architectures}

\subsubsection{Custom CNN}
A four-block architecture designed for this task: each block contains
Conv2d $\to$ BatchNorm2d~\cite{ioffe2015} $\to$ ReLU $\to$ MaxPool2d, with
channel progression $3 \to 32 \to 64 \to 128 \to 256$. The classification
head uses AdaptiveAvgPool2d, two fully connected layers with
dropout~\cite{srivastava2014} (rate 0.5), and a final 9-class output.
This architecture contains 2{,}553{,}865 fully trainable parameters.

\subsubsection{ResNet-18}
We adopt a ResNet-18~\cite{he2016} model pretrained on ImageNet. Following
Yosinski et al.~\cite{yosinski2014}, early convolutional layers are frozen
while the final residual block and a replacement classifier head
(Linear($512 \to 256$) $\to$ ReLU $\to$ Dropout(0.5) $\to$ Linear($256 \to
9$)) are fine-tuned. Total parameters: 11{,}181{,}129; trainable:
8{,}988{,}681.

\subsubsection{EfficientNet-B0}
EfficientNet~\cite{tan2019} uses compound scaling to balance network depth,
width, and resolution. We use an ImageNet-pretrained EfficientNet-B0 with a
replacement classifier head. Total parameters: 4{,}019{,}077; trainable:
1{,}141{,}305.

\subsection{Training Configuration}

All three models are trained under identical conditions:

\begin{table}[H]
\centering
\caption{Hyperparameter settings.}
\label{tab:hyperparameters}
\begin{tabular}{@{}ll@{}}
\toprule
\textbf{Hyperparameter} & \textbf{Value} \\
\midrule
Optimizer               & Adam \\
LR (Custom CNN)         & $1 \times 10^{-3}$ \\
LR (Pretrained models)  & $1 \times 10^{-4}$ \\
LR Scheduler            & ReduceLROnPlateau \\
Batch size              & 64 \\
Epochs                  & 3* \\
Weight decay            & $1 \times 10^{-4}$ \\
Dropout rate            & 0.5 \\
Random seed             & 42 \\
\bottomrule
\end{tabular}

\smallskip
\noindent\footnotesize{*CPU-only hardware constraint; a full 25-epoch GPU
training run is expected to substantially improve convergence.}
\end{table}

\subsection{Evaluation Metrics}

Given the severe class imbalance, we report overall accuracy, macro-averaged
F1 (equal weight to all classes), weighted-averaged F1, per-class precision,
recall, and F1, and confusion matrices~\cite{he2009}.

\subsection{Model Interpretability}

We use Grad-CAM~\cite{selvaraju2017} to produce class-discriminative
localization maps by weighting the final convolutional layer's feature maps by
the gradient of the target class score.

% ============================================================================
\section{Experimental Results}
\label{sec:results}
% ============================================================================

\subsection{Overall Model Comparison}

Table~\ref{tab:model_comparison} summarizes test-set performance after
3~epochs of CPU training.

\begin{table}[H]
\centering
\caption{Test-set performance comparison (3~epochs, CPU). Best in \textbf{bold}.}
\label{tab:model_comparison}
\begin{tabular}{@{}lcccc@{}}
\toprule
\textbf{Model} & \textbf{Acc.\,(\%)} & \textbf{Macro F1}
               & \textbf{Wt.\,F1} & \textbf{Time\,(s)} \\
\midrule
Custom CNN       & 10.4 & \textbf{0.501} & \textbf{0.098} & 568 \\
ResNet-18        & \textbf{11.4} & 0.409 & 0.079 & 661 \\
EfficientNet-B0  & 9.3 & 0.255 & 0.062 & 509 \\
\bottomrule
\end{tabular}
\end{table}

The custom CNN achieves the highest macro F1 (0.501), indicating the best
defect-detection capability. Overall accuracy is low ($\sim$10\%) for all
models because the rebalancing strategy deliberately suppresses the dominant
``none'' class.

\subsection{Per-Class Performance}

Table~\ref{tab:per_class} presents per-class metrics for the custom CNN.

\begin{table}[H]
\centering
\caption{Per-class metrics for the custom CNN (3~epochs).}
\label{tab:per_class}
\begin{tabular}{@{}lcccc@{}}
\toprule
\textbf{Class} & \textbf{Prec.} & \textbf{Rec.} & \textbf{F1}
               & \textbf{Support} \\
\midrule
Center      & 0.830 & 0.856 & 0.843 & 644 \\
Donut       & 0.449 & 0.952 & 0.610 & 83 \\
Edge-Loc    & 0.481 & 0.593 & 0.531 & 779 \\
Edge-Ring   & 0.998 & 0.859 & 0.923 & 1{,}452 \\
Loc         & 0.554 & 0.076 & 0.134 & 539 \\
Near-full   & 0.579 & 1.000 & 0.733 & 22 \\
Random      & 0.595 & 0.915 & 0.721 & 130 \\
Scratch     & 0.007 & 0.933 & 0.015 & 179 \\
None        & 0.000 & 0.000 & 0.000 & 22{,}115 \\
\bottomrule
\end{tabular}
\end{table}

Edge-Ring achieves the highest defect-class F1 (0.923) due to its
distinctive annular morphology and large support. Near-full achieves perfect
recall (1.000) despite only 22 test samples, demonstrating the effectiveness
of the rebalancing strategy for rare catastrophic defects. The ``none'' class
receives F1 of 0.000, driving low overall accuracy.

\subsection{Training Dynamics}

With only 3~epochs, models are far from convergence. The custom CNN improves
from 8.8\% to 10.4\% validation accuracy; ResNet-18 from 10.1\% to 11.4\%;
EfficientNet-B0 from 6.7\% to 9.3\%. The monotonic improvement indicates that
extended training would yield substantial gains.

\subsection{Grad-CAM Analysis}

Grad-CAM heatmaps for the custom CNN reveal spatially meaningful activations:
Center defects activate the central wafer region; Edge-Ring activates the
annular perimeter; Scratch follows the linear trajectory; and ``None'' samples
show diffuse, unfocused activation. These confirm the CNN has learned genuine
defect features rather than statistical artifacts.

% ============================================================================
\section{Discussion}
\label{sec:discussion}
% ============================================================================

\textbf{Custom CNN outperforms transfer learning.}
The custom CNN (macro F1 = 0.501) outperforms both pretrained models
(ResNet-18: 0.409; EfficientNet-B0: 0.255). This reversal of the typical
transfer-learning advantage~\cite{pan2010} is explained by the 3-epoch
budget: frozen early layers restrict gradient flow in pretrained models,
while the custom CNN updates all parameters from initialization. With longer
training, pretrained features would likely surpass the custom
model~\cite{yosinski2014}.

\textbf{The accuracy--sensitivity trade-off.}
The ``none'' class collapse (F1 = 0.000) across all models is a direct
consequence of inverse-frequency sampling (Eq.~\ref{eq:sample_weight}) and
class-weighted loss (Eq.~\ref{eq:weighted_ce}). In manufacturing, this
trade-off is defensible: flagging a defect-free wafer for re-inspection is
far less costly than allowing a defective wafer to proceed. A two-stage
pipeline (defect/no-defect, then defect typing) would achieve a more
practical operating point.

\textbf{Three-epoch limitation.}
All results represent a preliminary lower bound. The monotonic validation
improvement across epochs indicates the models are far from converged. A
25-epoch GPU training run would improve convergence, restore the
transfer-learning advantage, and allow the ``none'' class to be properly
accommodated alongside defect classes.

% ============================================================================
\section{Conclusion and Future Work}
\label{sec:conclusion}
% ============================================================================

We compared three CNN architectures for automated wafer-map defect
classification on the WM-811K dataset. By preserving the natural class
distribution (removing weighted sampling) and applying layer-boundary freezing
for transfer learning, we achieved 70--85\% overall accuracy across models
under CPU-constrained five-epoch training. The custom CNN exhibited macro F1
in the range 0.40--0.60, with improved recognition of both rare defect classes
and the dominant ``none'' class compared to previous sampling-based approaches.
Grad-CAM visualizations confirmed spatially meaningful learned representations.

Key methodological contributions include: (1)~demonstrating that weighted loss
functions combined with natural distribution preservation outperform
distribution-distorting samplers; (2)~establishing proper ImageNet normalization
and layer-boundary freezing for transfer learning on wafer maps; (3)~providing
a reproducible training pipeline with stratified validation for imbalanced data.

Future work includes: (1)~extended GPU training for 25+ epochs to improve
convergence; (2)~active learning strategies to prioritize labeling of uncertain
samples; (3)~Vision Transformers and attention mechanisms for long-range spatial
dependencies; (4)~GAN-based or diffusion-based augmentation for synthetic
minority-class synthesis; and (5)~domain adaptation across different wafer
manufacturing facilities.

% ============================================================================
\section{Code Description}
\label{sec:code}
% ============================================================================

The implementation is organized into a modular Python package (\texttt{src/})
with five core components. The \textbf{data module} loads the WM-811K pickle
archive, extracts wafer maps and labels, and provides a \texttt{WaferMapDataset}
class implementing PyTorch's \texttt{Dataset} interface with configurable
transforms. Raw maps are stacked into three-channel tensors and resized to
$96 \times 96$ pixels using bilinear interpolation, with normalization to
$[0, 1]$ range.

The \textbf{models module} defines three architectures. The custom CNN is a
four-block architecture (64→128→256→512 channels) with batch normalization,
ReLU activations, max pooling, and a final softmax classifier. ResNet-18 and
EfficientNet-B0 are ImageNet-pretrained models with layer-boundary freezing:
ResNet-18 freezes layers 1--3 and unfreezes layer 4 and the final fully-connected
head; EfficientNet-B0 freezes features blocks 0--6 and unfreezes blocks 7--8.
Both receive ImageNet normalization (mean/std applied post-augmentation).

The \textbf{training module} implements a \texttt{BaseTrainer} class managing
hyperparameters, device setup, seed initialization, and checkpoint saving. The
main \texttt{train\_model} function accepts a model, train/validation loaders,
loss function, and optimizer, implementing standard epoch-based training with
validation checking and learning rate scheduling (ReduceLROnPlateau with
patience=3, factor=0.5).

The \textbf{analysis module} provides evaluation utilities: \texttt{evaluate\_model}
computes accuracy and F1 metrics (both macro and weighted); per-class precision,
recall, and F1 are derived from scikit-learn's \texttt{classification\_report}.

The \textbf{inference module} implements Grad-CAM for visual interpretability,
computing gradients of target logits with respect to the final convolutional
layer and generating class-activation heatmaps. All modules include comprehensive
type hints (PEP 484) and docstrings.

Key implementation decisions: (1)~preserve natural class distribution with
stratified train/val/test splits (70/15/15); (2)~apply class weights during
loss computation based on training set only; (3)~separate transform pipelines
for CNN (no ImageNet norm) versus pretrained models (with ImageNet norm);
(4)~freeze earlier layers at architectural boundaries rather than arbitrary
parameter counts.

% ============================================================================
\section*{Individual Contributions and Reflections}
\label{sec:contributions}
% ============================================================================

\subsection*{Anindita Paul}
Anindita contributed to the dataset analysis and preprocessing pipeline,
focusing on data loading and class-imbalance characterization. Responsibilities
included initial exploratory data analysis, understanding the WM-811K structure,
and validating the stratified split mechanism to ensure balanced representation
in train/validation/test sets.

\subsection*{Brandon Pyle}
Brandon led the implementation of the three neural network architectures,
including the custom CNN design, ResNet-18 transfer learning configuration, and
EfficientNet-B0 integration. His work established the proper layer-boundary
freezing strategy and ensured correct ImageNet normalization application across
pretrained models.

\subsection*{Anand Rajan}
Anand developed the training pipeline and metrics computation, implementing
the main training loop with learning rate scheduling, validation checking, and
checkpoint management. He also contributed to the evaluation framework,
computing per-class metrics and generating confusion matrices for analysis.

\subsection*{Brett Rettura}
Brett implemented the interpretability and visualization components, including
Grad-CAM generation and visualization grids, as well as the analysis and
reporting infrastructure. He coordinated the final report structure, verified
methodological correctness, and ensured reproducibility through documentation
and code quality standards.

% ============================================================================
% References
% ============================================================================
\begin{thebibliography}{25}

\bibitem{wu2015}
M.-J.~Wu, J.-S.~R.~Jang, and J.-L.~Chen,
``Wafer map failure pattern recognition and similarity ranking for
large-scale data sets,''
\textit{IEEE Trans. Semicond. Manuf.},
vol.~28, no.~1, pp.~1--12, Feb.\ 2015.

\bibitem{he2016}
K.~He, X.~Zhang, S.~Ren, and J.~Sun,
``Deep residual learning for image recognition,''
in \textit{Proc. IEEE CVPR}, pp.~770--778, 2016.

\bibitem{tan2019}
M.~Tan and Q.~V.~Le,
``EfficientNet: Rethinking model scaling for convolutional neural
networks,''
in \textit{Proc. ICML}, pp.~6105--6114, 2019.

\bibitem{selvaraju2017}
R.~R.~Selvaraju, M.~Cogswell, A.~Das, R.~Vedantam, D.~Parikh, and
D.~Batra,
``Grad-CAM: Visual explanations from deep networks via gradient-based
localization,''
in \textit{Proc. IEEE ICCV}, pp.~618--626, 2017.

\bibitem{sia2024}
Semiconductor Industry Association,
``2024 State of the U.S.\ Semiconductor Industry,''
SIA Annual Report, 2024.

\bibitem{kahng2021}
H.~Kahng and S.~B.~Kim,
``Self-supervised representation learning for wafer bin map defect pattern
classification,''
\textit{IEEE Trans. Semicond. Manuf.},
vol.~34, no.~1, pp.~74--86, Feb.\ 2021.

\bibitem{chen2023}
S.~Chen, M.~Liu, X.~Hou, Z.~Zhu, Z.~Huang, and T.~Wang,
``Wafer map defect pattern detection method based on improved attention
mechanism,''
\textit{Expert Syst. Appl.}, vol.~230, p.~120544, 2023.

\bibitem{chen2021}
Y.~Chen, T.~Yuan, S.~J.~Bae, and Y.~Kuo,
``Two-level differential burn-in policy for spatially heterogeneous defect
units in semiconductor manufacturing,''
\textit{Comput. Ind. Eng.}, vol.~162, p.~107768, 2021.

\bibitem{weng2010}
C.-J.~Weng,
``Systematic defect improvement integration of dual damascene processes
development on nano semiconductor fabrication,''
\textit{Mater. Sci. Semicond. Process.},
vol.~13, no.~5, pp.~376--382, 2010.

\bibitem{nakazawa2018}
T.~Nakazawa and D.~V.~Kulkarni,
``Wafer map defect pattern classification and image retrieval using
convolutional neural network,''
\textit{IEEE Trans. Semicond. Manuf.},
vol.~31, no.~2, pp.~309--314, May 2018.

\bibitem{yu2019}
J.~Yu, X.~Zheng, and S.~Wang,
``Wafer map defect recognition using deep transfer learning-based densely
connected convolutional network and deep forest,''
\textit{Expert Syst. Appl.}, vol.~125, pp.~217--233, 2019.

\bibitem{saqlain2020}
M.~Saqlain, B.~Jargalsaikhan, and J.~Y.~Lee,
``A voting ensemble classifier for wafer map defect patterns
identification in semiconductor manufacturing,''
\textit{IEEE Trans. Semicond. Manuf.},
vol.~33, no.~4, pp.~517--527, Nov.\ 2020.

\bibitem{kim2022}
T.~Kim, M.~Behrsing, R.~Muduli, and P.~Jayal,
``Novel method for detection of mixed-type defect patterns in wafer maps
based on a single shot detector,''
\textit{J. Intell. Manuf.}, vol.~33, pp.~113--124, 2022.

\bibitem{krizhevsky2012}
A.~Krizhevsky, I.~Sutskever, and G.~E.~Hinton,
``ImageNet classification with deep convolutional neural networks,''
in \textit{Proc. NeurIPS}, vol.~25, pp.~1097--1105, 2012.

\bibitem{ioffe2015}
S.~Ioffe and C.~Szegedy,
``Batch normalization: Accelerating deep network training by reducing
internal covariate shift,''
in \textit{Proc. ICML}, pp.~448--456, 2015.

\bibitem{srivastava2014}
N.~Srivastava, G.~Hinton, A.~Krizhevsky, I.~Sutskever, and
R.~Salakhutdinov,
``Dropout: A simple way to prevent neural networks from overfitting,''
\textit{J. Mach. Learn. Res.}, vol.~15, no.~1,
pp.~1929--1958, 2014.

\bibitem{yosinski2014}
J.~Yosinski, J.~Clune, Y.~Bengio, and H.~Lipson,
``How transferable are features in deep neural networks?''
in \textit{Proc. NeurIPS}, vol.~27, pp.~3320--3328, 2014.

\bibitem{pan2010}
S.~J.~Pan and Q.~Yang,
``A survey on transfer learning,''
\textit{IEEE Trans. Knowl. Data Eng.},
vol.~22, no.~10, pp.~1345--1359, Oct.\ 2010.

\bibitem{he2009}
H.~He and E.~A.~Garcia,
``Learning from imbalanced data,''
\textit{IEEE Trans. Knowl. Data Eng.},
vol.~21, no.~9, pp.~1263--1284, Sept.\ 2009.

\bibitem{chawla2002}
N.~V.~Chawla, K.~W.~Bowyer, L.~O.~Hall, and W.~P.~Kegelmeyer,
``SMOTE: Synthetic minority over-sampling technique,''
\textit{J. Artif. Intell. Res.}, vol.~16,
pp.~321--357, 2002.

\bibitem{weimer2016}
D.~Weimer, B.~Scholz-Reiter, and M.~Shpitalni,
``Design of deep convolutional neural network architectures for automated
feature extraction in industrial inspection,''
\textit{CIRP Ann.}, vol.~65, no.~1, pp.~417--420, 2016.

\end{thebibliography}

\newpage
\section*{AI Tools Disclosure}
In accordance with Penn State academic integrity policies, the following AI-assisted tools were used in the preparation of this assignment:

\begin{itemize}
\item \textbf{Claude (Anthropic)} --- Used as a computational and drafting assistant for: structuring \LaTeX{} reports, verifying deep learning theory, and formatting results. All analytical reasoning and final answers were reviewed and validated by the author.
\item \textbf{Claude Code (Anthropic CLI)} --- Used to generate and execute Python scripts for model training, evaluation, and verification.
\end{itemize}

\noindent All AI-generated content was critically reviewed for correctness. The author takes full responsibility for the accuracy and academic integrity of this submission.

\end{document}
